% ============================================================
% §4 · Instalación y configuración — Manual de Usuario K-QGMIP
% ============================================================

\section{Instalación y configuración}
\label{sec:instalacion}

K-QGMIP se distribuye como \textbf{código fuente} a través de un repositorio Git.
No existe un instalador ejecutable: el usuario descarga el repositorio, crea un entorno
virtual de Python y ejecuta la instalación de dependencias con un único comando.

% ------------------------------------------------------------
\subsection{Paso 1 — Obtener el código fuente}
% ------------------------------------------------------------

Clone el repositorio con Git o descargue el archivo ZIP desde la plataforma de distribución
del proyecto:

\begin{lstlisting}[language=bash]
git clone https://github.com/CristianDavGon/algoritmos_final1.git

cd "Proyecto V03 FINAL"
\end{lstlisting}

Si no tiene Git instalado, puede descargar el ZIP directamente desde la interfaz web del
repositorio y descomprimirlo en la ruta de su preferencia.

% ------------------------------------------------------------
\subsection{Paso 2 — Verificar Python 3.12}
% ------------------------------------------------------------

Confirme que tiene Python 3.12 o superior instalado:

\begin{lstlisting}[language=bash]
python --version
\end{lstlisting}

La salida debe mostrar \texttt{Python 3.12.x} o superior. Si no cumple este requisito,
descargue e instale la versión correcta desde \texttt{python.org} antes de continuar.

% ------------------------------------------------------------
\subsection{Paso 3 — Crear el entorno virtual}
% ------------------------------------------------------------

Se recomienda usar \textbf{\texttt{uv}}, el gestor de entornos del proyecto, para aislar
las dependencias y evitar conflictos con otros proyectos Python en el sistema.

\textbf{Con \texttt{uv} (recomendado):}

\begin{lstlisting}[language=bash]
pip install uv
cd code
uv venv
\end{lstlisting}

\textbf{Con \texttt{venv} estándar (alternativa):}

\begin{lstlisting}[language=bash]
cd code
python -m venv .venv
\end{lstlisting}

Esto crea una carpeta \texttt{.venv/} dentro de \texttt{code/} con un intérprete Python
aislado. No modifique ni elimine esta carpeta manualmente.

% ------------------------------------------------------------
\subsection{Paso 4 — Activar el entorno virtual}
% ------------------------------------------------------------

Antes de instalar dependencias o ejecutar cualquier script, active el entorno:

\textbf{Windows (PowerShell):}
\begin{lstlisting}[language=bash]
.venv\Scripts\Activate.ps1
\end{lstlisting}

\textbf{Windows (símbolo del sistema):}
\begin{lstlisting}[language=bash]
.venv\Scripts\activate.bat
\end{lstlisting}

\textbf{Linux / macOS:}
\begin{lstlisting}[language=bash]
source .venv/bin/activate
\end{lstlisting}

El prefijo \texttt{(.venv)} en el prompt del terminal confirma que el entorno está activo.
Recuerde activarlo cada vez que abra una nueva sesión de terminal.

% ------------------------------------------------------------
\subsection{Paso 5 — Instalar las dependencias}
% ------------------------------------------------------------

Con el entorno activo, instale todas las dependencias declaradas en \texttt{pyproject.toml}:

\textbf{Con \texttt{uv} (recomendado):}
\begin{lstlisting}[language=bash]
uv sync
\end{lstlisting}

\textbf{Con \texttt{pip} (alternativa):}
\begin{lstlisting}[language=bash]
pip install -e .
\end{lstlisting}

\texttt{uv sync} resuelve las versiones exactas desde el archivo \texttt{uv.lock},
garantizando reproducibilidad. \texttt{pip install -e .} instala en modo editable y
resuelve versiones mínimas desde \texttt{pyproject.toml}.

% ------------------------------------------------------------
\subsection{Paso 6 — Verificar la instalación}
% ------------------------------------------------------------

Compruebe que las bibliotecas principales están disponibles ejecutando desde la carpeta
\texttt{code/}:

\begin{lstlisting}[language=Python]
python -c "import numpy, scipy, pandas, pyphi; print('OK')"
\end{lstlisting}

Si la salida es \texttt{OK}, el entorno está listo para ejecutar K-QGMIP.
Si algún paquete falla, repita el Paso 5 asegurándose de que el entorno virtual esté activo.

% ------------------------------------------------------------
\subsection{Nota sobre paralelización}
% ------------------------------------------------------------

K-QGMIP corre íntegramente sobre \textbf{CPU} sin configuración adicional. Si la máquina
dispone de múltiples núcleos, el \textit{framework} puede aprovecharlos para paralelizar
la evaluación de particiones candidatas; en ese caso no se requiere ningún hardware
especializado (GPU) ni ajustes manuales: el sistema detecta automáticamente los núcleos
disponibles y los utiliza durante la ejecución.
